\documentclass[10pt]{article}
\usepackage[T1]{fontenc}
\usepackage[utf8]{inputenc}
\usepackage[french]{babel}
\usepackage{amsmath,amssymb,mathrsfs,mathpazo}
\usepackage{geometry}
\geometry{a4paper,margin=2cm}
\usepackage{graphicx,array,tabularx,tabularray,colortbl,multirow,multicol,booktabs,enumitem,pifont,ulem,fancybox,diagbox,cancel}
\usepackage{pstricks,pstricks-add,pst-plot,pst-tree,pst-node,pst-text,pst-func,pst-math,pst-bspline,pst-3dplot,pst-eucl,pst-solides3d}
\let\pscalculate\undefined
\usepackage{pst-bezier}
\usepackage{tikz,pgfplots,tkz-tab}
\usetikzlibrary{arrows,arrows.meta,patterns,calc,positioning,trees,decorations.markings,shapes.geometric,shapes.misc,shapes.arrows,intersections,decorations.pathmorphing,decorations.pathreplacing,backgrounds,shadows,babel}
\pgfplotsset{compat=1.18}
\usepgfplotslibrary{fillbetween}
\usepackage[np]{numprint}
\usepackage{listings,makecell,pas-tableur,twoopt}
\let\barre\undefined
\usepackage[pstricks]{tablvar}
\definecolor{darkgreen}{RGB}{0,100,0}
\providecommand{\red}{\color{red}}
\providecommand{\blue}{\color{blue}}
\providecommand{\green}{\color{green}}
\providecommand{\black}{\color{black}}
\providecommand{\white}{\color{white}}
\providecommand{\gray}{\color{gray}}
\providecommand{\cyan}{\color{cyan}}
\providecommand{\magenta}{\color{magenta}}
\providecommand{\yellow}{\color{yellow}}
\providecommand{\violet}{\color{violet}}
\providecommand{\orange}{\color{orange}}
\providecommand{\euro}{\texteuro}
\providecommand{\EUR}{\texteuro}
\providecommand{\ieme}{\textsuperscript{e}}
\providecommand{\degres}{\ensuremath{^\circ}}
\providecommand{\degree}{\ensuremath{^\circ}}
\providecommand{\vv}[1]{\overrightarrow{#1}}
\providecommand{\Coord}[1]{#1}
\providecommand{\decosix}{}
\providecommand{\starredbullet}{\ensuremath{\bullet}}
\pagestyle{empty}
\setlength{\parindent}{0pt}
\usepackage{adjustbox}\DeclareUnicodeCharacter{2236}{:}\DeclareUnicodeCharacter{2190}{\ensuremath{\leftarrow}}
\providecommand{\R}{\mathbb{R}}
\providecommand{\N}{\mathbb{N}}
\providecommand{\Z}{\mathbb{Z}}
\providecommand{\Q}{\mathbb{Q}}
\providecommand{\D}{\mathbb{D}}
\providecommand{\C}{\mathbb{C}}
\providecommand{\e}{{\mathrm{e}}}
\providecommand{\vect}[1]{\overrightarrow{#1}}
\providecommand{\vectt}[1]{\overrightarrow{\mathrm{#1}}}
\providecommand{\vvt}[1]{\overrightarrow{\mathrm{#1}}}
\providecommand{\np}[2][]{{#2~#1}}
\providecommand{\vv}[1]{{\overrightarrow{#1}}}
\renewcommand{\d}{{\mathrm{d}}}
\providecommand{\degres}{{^\circ}}
\providecommand{\degre}{{^\circ}}
\providecommand{\blue}{\color{blue}}
\providecommand{\red}{\color{red}}
\providecommand{\enskip}{\quad}
\providecommand{\llbracket}{\lbrack\!\lbrack}
\providecommand{\rrbracket}{\rbrack\!\rbrack}
\providecommand{\euro}{{\text{€}}}
\providecommand{\no}{{\text{n°}}}
\providecommand{\interval}[2]{{[#1;#2]}}
\providecommand{\intervFO}[2]{{[#1;#2[}}
\providecommand{\dotfill}{\ldots}

\providecommand{\ds}{\displaystyle}
\providecommand{\textding}[1]{#1}
\providecommand{\Oij}{\(\mathrm{O};\vec{\imath},\vec{\jmath}\)}
\providecommand{\Oijk}{\(\mathrm{O};\vec{\imath},\vec{\jmath},\vec{k}\)}
\providecommand{\Ouv}{\(\mathrm{O};\vec{u},\vec{v}\)}
\providecommand{\intervFF}[2]{\ensuremath{[#1;#2]}}
\renewcommand{\euro}{\ifmmode\text{\texteuro}\else\texteuro\fi}
\DeclareUnicodeCharacter{00B0}{\ifmmode^{\circ}\else\textdegree\fi}
\DeclareUnicodeCharacter{20AC}{\ifmmode\text{\texteuro}\else\texteuro\fi}
% Énoncés : Éducation nationale. Transcription et corrigés : APMEP, Denis Vergès et contributeurs. Découpage et classement : CoopMaths, Sébastien Lozano, Éric Elter et contributeurs. Adaptation de la présentation : Maths Explorer.
% CC BY-SA : https://coopmaths.fr/www/about/
% https://coopmaths.fr/alea/static/e3c/2020/tex/e3c_2020_00_sujet27_2.tex
\begin{document}
\noindent\textit{APMEP / CoopMaths. CC BY-SA. Présentation adaptée par Maths Explorer.}\par\medskip
\medskip

Dans le plan muni d'un repère, on a tracé la courbe représentative $\mathcal{C}_f$ d'une fonction $f$ définie et dérivable sur $\R$. On note $f'$ la dérivée de $f$. 

On sait que la courbe $\mathcal{C}_f$ admet exactement deux tangentes horizontales :

\begin{itemize}
\item l'axe des abscisses comme tangente à la courbe $\mathcal{C}_f$ au point $A(-1~;~0)$ ;
\item la droite $T_B$ comme tangente à la courbe $\mathcal{C}_f$ au point $B\left(\dfrac{1}{3}~;~-\dfrac{32}{27}\right)$.
\end{itemize}

\begin{center}
\psset{xunit=2.5cm,yunit=1.5cm,labelFontSize=\scriptstyle,showorigin=false,comma=true}
\begin{pspicture}(-3,-3)(3,2.6)
\multido{\n=-2.5+0.1}{50}{\psline[linewidth=0.35pt,linecolor=lightgray](\n,-2.9)(\n,2.4)}
\multido{\n=-2.9+0.1}{54}{\psline[linewidth=0.35pt,linecolor=lightgray](-2.5,\n)(2.4,\n)}
\multido{\n=-2.5+0.5}{10}{\psline[linewidth=0.45pt](\n,-2.9)(\n,2.4)}
\multido{\n=-2.5+0.5}{10}{\psline[linewidth=0.45pt](-2.5,\n)(2.4,\n)}
\psaxes[linewidth=0.95pt,Dx=0.5,Dy=0.5]{->}(0,0)(-2.5,-2.9)(2.5,2.4)
\def\Func{x x x 1 add mul 1 sub mul 1 sub }
\psplot[plotpoints=2000,linewidth=1.25pt,linecolor=red]{-2}{1.4}{\Func}
\psline[linewidth=1.25pt](-2.5,-1.185185)(2.4,-1.185185)
\psdots[dotstyle=Mul,dotscale=2.06,linecolor=blue](-1,0)(0.333333,-1.185185)
\uput[u](-1,0){$A$}\uput[d](0.333333,-1.185185){$B$}\uput[dl](0,0){O}\uput[r](1.01,1.3){\red $\mathcal{C}_f$}
\uput[d](2.2,-1.25){$T_B$}
\end{pspicture}
\end{center}

\begin{enumerate}
\item  Par lecture graphique, donner les solutions de l'équation $f(x)=0$.
\end{enumerate}

La fonction $f$ est définie sur $\R$ par $f(x)=x^3+x^2-x-1$. On note $f'$ la dérivée de $f$.

\begin{enumerate}[resume]
\item Déterminer $f'(x)$ pour tout réel $x$.
\item En déduire le tableau de variations de $f$.
\item En utilisant ce qui précède, déterminer la position relative de la courbe $\mathcal{C}_g$ de la fonction $g$ définie sur $\R$ par $g(x)=x^3+x^2$ et de la droite $D$ d'équation $y=x+1$.
\end{enumerate}

\vspace{0,5cm}

\end{document}
